\documentclass[../../../include/open-logic-section]{subfiles}

\begin{document}

\olfileid[pt]{sth}{z}{milestone}
\olsection{$\Z^-$: um marco}

Voltaremos a \stagesinf{} na parte seguinte. Contudo, com o
Axioma do Infinito, atingimos um marco importante. Temos agora
todos os axiomas necessários para a teoria $\Zminus$. Em detalhe:

\begin{defn}
A teoria $\Zminus$ tem estes axiomas: Extensionalidade, União, Pares, Conjuntos potência, Infinito, e todas as instâncias do esquema de Separação.
\end{defn}

O nome abrevia \emph{Zermelo} (\emph{menos} algo
que abordaremos mais tarde). Zermelo merece a honra, pois
formulou essencialmente esta teoria no seu trabalho de
\citeyear{Zermelo1908Untersuchungen}.\footnote{Para comentários interessantes
sobre a história e os pormenores técnicos, veja \citet[Apêndice
A]{Potter2004}.}

Esta teoria é bastante poderosa para permitir-nos fazer uma quantidade enorme de
matemática. Em particular, \emph{deverá} voltar atrás na
\olref[sfr][][]{part} e convencer-se de que tudo o que fizemos,
ingenuamente, pode ser feito de modo mais formal dentro de~$\Zminus$. (Depois de o fazer durante algum tempo, poderá querer avançar e ler a
\olref[sth][z][arbintersections]{sec}.) Por isso, daqui em diante,
e sem mais comentários, assumiremos estar a trabalhar
em $\Zminus$ (pelo menos).

\end{document}
